\documentclass[12pt]{extarticle}

\usepackage[english]{babel}
\usepackage[utf8x]{inputenc}
\usepackage{amsmath}
\usepackage{graphicx}
\usepackage{lmodern}
\usepackage[T1]{fontenc}

\usepackage[colorinlistoftodos]{todonotes}
\usepackage[margin=1.2in]{geometry}
\usepackage{textcomp}
\newcommand{\textapprox}{\raisebox{0.5ex}{\texttildelow}}

\usepackage{multirow}
\usepackage{amsmath,amssymb}
\usepackage{centernot}

\usepackage{minted}


\title{CSCI-304-02 Computer Organization: Assignment 3}

\begin{document}
\maketitle


\textbf{Due: 11:59pm 10/20/2023}

\textbf{Submission: Please submit a .c file (C source code), a .txt file (explain compile details) and a PDF file (your answers to the questions) on Blackboard.}\\


Show your work for full credit. All submission must be completely your own work.\\
\\

1. [50 points] \textbf{Project: Simple Cipher in C}

In this project, you are tasked with creating a C program to encode a provided message. The detailed requirements are as follows:

\begin{itemize}
    \item \textbf{Encoding Scheme:}
    \begin{itemize}
        \item Reverse each ``word'' in the given phrase. A ``word'' is defined as a sequence of alphabetic characters separated by one or more non-alphabetic characters.
        \item Maintain the position of all non-alphabetic characters. As a result, the length of the output string should match the length of the input string.
    \end{itemize}
    
    \item \textbf{Handling Underscores:}
    \begin{itemize}
        \item The input message will use underscores (`\_') instead of spaces to facilitate reading it as a single text value.
        \item In your encoded output, convert all underscores back to spaces.
    \end{itemize}
    
    \item \textbf{Case Handling:}
    \begin{itemize}
        \item Ensure all alphabetic characters in the output are in uppercase.
    \end{itemize}
    \item \textbf{Program Termination:}
    \begin{itemize}
        \item Each input should end with a newline.
        \item Users can exit the program by entering the word ``quit''.
    \end{itemize}
\end{itemize}

\clearpage

Example inputs -> outputs:
\begin{itemize}
    \item War\_eagle -> RAW ELGAE
    \item Reading\_records\_of\_variable\_length? -> GNIDAER SDROCER FO ELBAIRAV HTGNEL? 
    \item Have\_fun.\_Doing\_this\_lab\_:) -> EVAH NUF. GNIOD SIHT BAL :)
    \item Words\%end*with\^{}non-alpha...characters!!! -> \\SDROW\%DNE*HTIW\^{}NON-AHPLA...SRETCARAHC!!!
\end{itemize}


CONSTRAINTS:
\begin{itemize}
    \item  The length of the input message to be encoded will not exceed 100 characters, including the new line character.
    \item Use pointers to manipulate the location of each character in the array(s)
    \item All code should be submitted in a single source file named \textit{simple\_cipher.c}. Please specify the command you use to compile your program either in a \textit{.txt file} or by submitting a \textit{Makefile} along with your solution.
\end{itemize}

Your program can be compiled with the command:\\
\begin{minted}{bash}
  $ gcc simple_cipher.c –o simple_cipher
\end{minted}
..and run with the command to test with the sample input file on the blackboard:
\begin{minted}{bash}
  $ ./simple_cipher < test2.txt
\end{minted}
      
      
or with the command (to save the output):
\begin{minted}{bash}
  $ ./simple_cipher < text2.txt > mytestout2.txt
\end{minted}
      
Compare results with sample results on the class webpage using the command:
\begin{minted}{bash}
  $ diff mytestout2.txt testout2.txt
\end{minted}
\pagebreak

2. [6 points] Consider the following C code. What is the value of variable $s$ after the for loop? Explain how you get the value of $s$.  \\%Use repeated division for the initial conversion from decimal to binary

\begin{minted}{c}
int i;  /* 32 bit integer */
unsigned int s=1; /* 32 bit unsigned integer */
for (i = -1 ; i < 0; i--) 
{
    s = s+1; 
}
\end{minted}
$i$ is a 32 bit signed integer, which can be from $-(2^{31})$ to $(2^{31}-1)$ \\ (-2,147,483,648 to 2,147,483,647)\\
$s$ is a 32 bit unsigned integer, which can be from $0$ to $2^{32}-1$ (0 to 4,294,967,296)\\~\\
$s$ starts as 1, and $i$ starts as -1\\
The loop will run from $i=-1$ to $i=-2,147,483,648$, terminating when $i$ overflows to zero\\
Therefore, the loop runs, incrementing $s$, 2,147,483,648 times\\
At the end, s will be 2,147,483,649 (1+ 2,147,483,648)\\

\pagebreak
3. [12 points] Show how values are stored in memory using the Sun, IA32, and x86-64 formats, as demonstrated on slide 80 of \texttt{Lec04\_integers\_02072023}. Assume each value is declared as a \texttt{long int} in C. Both Sun and IA32 use 32-bit encoding for \texttt{long int}, while x86-64 uses 64-bit encoding.\\

a. -862 \\
-862$_{10}$ = 1100 1010 0010$_{2}$ = (-2048 + 1024 + 128 + 32 + 2)\\
1100 1010 0010$_{2}$ = CA2$_{16}$\\
expand to 32 bits for IA32 and Sun: FFFFFCA2$_{16}$, \\
and to 64 bits for x86-64: FFFFFFFF FFFFFCA2$_{16}$\\

b. 3855 \\
3855$_{10}$ = 1111 0000 1111$_{2}$ (2048+1024+512+256+8+4+2+1)\\
1111 0000 1111$_{2}$ = F0F$_{16}$ \\
expand to 32 bits for IA32 and Sun: 00000F0F$_{16}$, \\
and to 64 bits for x86-64: 00000000 00000F0F$_{16}$\\

c. -6691 \\
-6691$_{10}$ = 10 0101 1101 1101$_{2}$ (-8192+(1024+256+128+64+16+8+4+1))\\
(Expand to 16 bit before conversion): 1110 0101 1101 1101$_{2}$ = E5DD$_{16}$ \\
expand to 32 bits for IA32 and Sun: FFFFE5DD$_{16}$, \\
and to 64 bits for x86-64: FFFFFFFF FFFFE5DD$_{16}$\\

\begin{center}
    \begin{tabular}{|c|c|c||c|c|c||c|c|c|}
    \hline
        \multicolumn{3}{|c||}{-862} &  \multicolumn{3}{|c||}{3855}  &  
        \multicolumn{3}{|c|}{-6691}  \\
    \hline
        IA32 & x86-64 & Sun & IA32 & x86-64 & Sun & IA32 & x86-64 & Sun\\
    \hline \hline
        A2 & A2 & FF & 0F & 0F & 00 & DD & DD & FF\\
        FC & FC & FF & 0F & 0F & 00 & E5 & E5 & FF\\
        FF & FF & FC & 00 & 00 & 0F & FF & FF & E5\\
        FF & FF & A2 & 00 & 00 & 0F & FF & FF & DD\\
           & FF &    &    & 00 &    &  & FF & \\
           & FF &    &    & 00 &    &  & FF & \\
           & FF &    &    & 00 &    &  & FF & \\
    \hline
    \end{tabular}\\
    IA32, x86-64, and Sun representations
\end{center}


\pagebreak



4. [6 points] Translate the following unsigned binary numbers into their base 10 equivalents. Use fractions in lowest terms to represent values on the right-hand side of the binary point.\\
\\
a. 0.0101011 $ = \frac{43}{128}$\\
$\frac{1}{4} +\frac{1}{16} +\frac{1}{64} + \frac{1}{128} = \frac{32+8+2+1}{128} = \frac{43}{128}$
\\
\\
b. 11100.0101 $= 28 \frac{5}{16}$\\
$16 + 8 + 4 + \frac{1}{4} + \frac{1}{16} = 28 \frac{5}{16}$
\\
\\


\pagebreak

5. [6 points] Translate the following decimal numbers into their binary equivalents using a binary point.
\\
\\
a. $79 \frac{1}{2}$\\
$79 \frac{1}{2} = 64 + 8 + 4 + 2 + 1 + \frac{1}{2} = $ 1001111.1$_2$
\\
\\
b. $62 \frac{9}{16}$\\
$62 \frac{9}{16} = 32 + 16 + 8 + 4 + 2 + \frac{1}{2} + \frac{1}{16} =$ 111110.1001$_2$
\\
\\


\pagebreak

6. [20 points] Consider the following C code. Assume that variables are stored in memory starting at location 1000, in order and with no gaps between variables. Assume all int’s and pointers consume 4 bytes each.

a. Show the memory layout and values for the variables at position A.

b. Show the memory layout and values for the variables at position B.

c. Show the memory layout and values for the variables at position C.

d. What is printed by the program?
\begin{minted}{c}
#include <stdio.h>
void main ()
{
    int x = 3, y = 0;
    int *p = &x;
    int *q = &y;
    int a[3] = {-1, 2, 6};
    /* position A */
    *q += 1;
    q--;
    *q += 2;
    /* position B */
    q = &a[1];
    *q = 7;
    p = q++;
    *p = 4;
    *q = 5;
    /* position C */
    printf ("%d, %d, %d\n", &p, p, *p);
    printf ("%d, %d, %d\n", &q, q, *q);
}

\end{minted}
\textbf{a-c. }\\
\begin{tabular}{|c|c|c|c|c|}
    \hline
     variable & address & value at A & value at B & value at C  \\
     \hline \hline
     x & 1000 & 3 & 5 & 5  \\
     y & 1004 & 0 & 1 & 1  \\
     *p & 1008 & 1000 & 1000 & 1020  \\
     *q & 1012 & 1004 & 1000 & 1024  \\
     a[0] & 1016 & -1 & -1 & -1  \\
     a[1] & 1020 & 2 & 2 & 4  \\
     a[2] & 1024 & 6 & 6 & 5  \\
     \hline
\end{tabular}\\
\textbf{d. }Program prints: \\
1008, 1020, 4\\
1012, 1024, 5

\end{document}